% !TEX root = ../../main/main.tex
\hypearticle
{O Teste de Turing foi superado? O que as máquinas nos ensinam sobre inteligência}
{Caio do Nascimento Barros}
{

\subtitle{Uma Pergunta Trocada por um Jogo}

\hspace*{1.5em}Em 1950, Alan Turing abriu um célebre artigo com a pergunta: máquinas podem pensar? Achando a questão vaga demais, já que ninguém sabe definir bem o que é pensar, ele a substituiu por um experimento, o jogo da imitação. Um juiz humano conversa apenas por texto com um computador e uma pessoa e tenta adivinhar quem é quem. Se a máquina for indistinguível de um humano, por que negar a ela o status de entidade pensante? Quem assistiu a \textit{Blade Runner}, em que o teste Voight-Kampff distingue humanos de replicantes por meio de perguntas, reconhece o clima.

\subtitle{Quando o Teste Vira Manchete}

\hspace*{1.5em}Turing propôs o teste como experimento mental e deixou detalhes em aberto. Arriscou apenas uma previsão: em cerca de 50 anos, um juiz médio acertaria no máximo 70\% das vezes após cinco minutos, ou seja, seria enganado em 30\%. Em 2014, na Royal Society, o chatbot Eugene Goostman, que fingia ser um adolescente ucraniano, enganou 10 de 30 juízes (33,3\%) e os organizadores declararam o teste superado. Especialistas discordaram: com tão pouco tempo e juízes de experiência desigual, o teste mediu a ingenuidade humana, não a inteligência da máquina. É o chamado efeito ELIZA, nome de um chatbot dos anos 1960 que, mesmo simplíssimo, convenceu muita gente de que era um psicoterapeuta compreensivo.

\hspace*{1.5em}Hoje, manchetes dizem que ChatGPT e Claude já passaram no teste, mas o significado da expressão se enfraqueceu. Em um estudo de 2024 com 500 participantes, em formato de dois jogadores, o GPT-4 foi tomado por humano em 54\% das rodadas de cinco minutos, contra 67\% dos humanos reais. Como o critério era enganar mais da metade dos juízes, ele \textit{passou}. Já o jogo original, com três participantes, juízes especialistas e mais tempo, continua sem vencedor.

\subtitle{Fluência Não É Inteligência}

\hspace*{1.5em}A história da IA ensina a desconfiar dessa intuição. Em 1958, pioneiros como Newell e Simon achavam que uma máquina de xadrez exigiria algo próximo da inteligência humana completa. Décadas depois, o Deep Blue venceu Garry Kasparov por força bruta, longe de qualquer inteligência geral. Há ainda evidências da neurociência: Ev Fedorenko, do MIT, e colaboradores mostram que as redes cerebrais da linguagem são em grande parte separadas das do raciocínio e do senso comum. Assumir que falar bem implica pensar bem seria, nas palavras deles, uma falácia. Isso não torna o tema inofensivo: usar IA generativa para se passar por humano, espalhando desinformação ou aplicando golpes, é um risco real.

\subtitle{O Que o Futuro Nos Reserva?}

\hspace*{1.5em}Turing previu que, no fim do século XX, falar em máquinas pensantes não causaria mais estranhamento. Pode ser que a previsão esteja errada por algumas décadas, ou que o que mude seja a nossa própria ideia de pensar. A segunda hipótese é a mais animadora: cada vez que uma máquina nos surpreende, somos forçados a refinar o que entendemos por inteligência, mais complexa e sutil do que Turing imaginava. Em \textit{Blade Runner}, os replicantes forçam os humanos a se perguntarem o que os torna humanos. Algo parecido acontece conosco: ao perguntar o que as máquinas conseguem, talvez aprendamos mais sobre nós mesmos.

}

\highlight{
    Para saber mais sobre o Teste de Turing:
}

\begin{hypeitemize}
\item \href{https://www.science.org/doi/10.1126/science.adq9356}{Ensaio original de Melanie Mitchell na \textit{Science}};
\item \href{https://aiguide.substack.com/p/the-turing-test-and-our-shifting}{Post da autora sobre o ensaio no AI Guide}.
\end{hypeitemize}

\originalsource{https://www.science.org/doi/10.1126/science.adq9356}